\section{担保物权}
\begin{enumerate}
  % https://zhuanlan.zhihu.com/p/591635305
  \item 要点。标的物（不同，担保规则不同）+ 顺序。
  \item 用他人之物或财产权利担保自己债权实现的限制物权。本质\zd{优先受偿权}。
  \item 担保（从）物权 \yjt 保障债权（主）的全面实现（保全）。必然 2 层法律关系。
  \item 担保人、债务人，债权人、担保权人合一。\al{按揭买房}。
        \begin{enumerate}
          \item 借款合同。主合同，债权效力 \yjt 请求权 \yjt 请求进行房屋抵押。
          \item 房屋抵押合同。从合同，债权效力 \yjt 请求办理登记手续。
          \item 房屋抵押登记。从合同 \yjt 从权力，债权 \yjt 物权（房屋抵押权）\yjt 优先受偿。
        \end{enumerate}
  \item \zd{担保的从属性}。发生、变更、转让、消灭、\zd{效力、范围（责任 <= 主债务人）}。
        \begin{enumerate}
          \item 强制性规定。不允许当事人做任何改动或创设，否则一律无效。
          \item 有担保无独立，有独立无担保。\zd{独立担保，违法无效}。例外：金融机构独立保函，不考。
          \item 主合同无效，担保合同也无效。
          \item 过度担保，权利滥用，违法无效。实践常见：\zd{针对担保责任约定专门的违约责任。另行支付；流押流置}。
        \end{enumerate}
  \item 不可分性。主债权尚未清偿完毕，担保物一个都不能少。
  \item \zd{物上代位权}。
        \begin{enumerate}
          \item 有代位物。\zd{保险金}（喜欢考）、赔偿金、补偿金等。
          \item 没有\blh。担保物毁损、灭失、被征收，担保物权消灭，\zd{从权利消灭，但主权利依然存在}（从属性）。
        \end{enumerate}
  \item 流押流置条款无效。约定标的物所有权归债权人所有。
  \item 担保合同无效的后果。合同相对性原理。
        \begin{enumerate}
          \item 不承担担保责任（物权效力）、违约责任/担保义务（债权效力）。
          \item 可能承担缔约过失责任；不能清偿的部分，补充责任。
        \end{enumerate}
  \item 动产担保 7 大家族。动产抵押、动产浮动抵押、超级动产抵押；\zd{动产质押；动产留置}；动产所有权保留买卖；动产融资租赁。
  \item 抵押权。必考，担保之王。\zd{不需要移转标的物的占有}，有抵押无交付，有交付无抵押，抵押和交付水火不容。
        \begin{enumerate}
          \item 抵押财产。法律不禁止；\zd{可交换性}，\al{公益法人的公益设施}。
          \item 不动产抵押权。\zd{登记生效}主义 。
                \begin{enumerate}
                  \item 不动产不允许对抵押财产进行概括性描述；动产允许 \yjt 动产浮动抵押。\al{抵押为房屋}。
                  \item 3 种：房、地、房和地（房地一体抵押）。
                  \item 新增建筑，登记时有啥算啥，新增的无优先受偿。
                \end{enumerate}
          \item 动产。有效合同生效时 = 动产抵押权 + \zd{登记 > 第三人}。
                \begin{enumerate}
                  \item 动产抵押权。不能拍卖，不能请求返还。
                  \item 动产浮动抵押/商事抵押权。核心点标的物 \yjt 现有的、\zd{即将有的}生产设备、原材料、材料、半成品。
                        \begin{enumerate}
                          \item 登记了也不能对抗\zd{正常经营买受人}（正常做生意，事出反常必有妖）。
                          \item 卖方 \yjt \zd{营业执照明确记载（不能是等）、持续销售}；买方 \yjt \zd{已支付合理对价；已取得}。
                        \end{enumerate}
                  \item 超级动产抵押权/购买价金担保权。\zd{买卖合同之债优先于借款合同之债}。
                        \begin{enumerate}
                          \item \zd{动产标的物来源于买卖合同的出卖人}，优先于其他债权人。
                          \item 交付标的后的 10 日内登记 \yjt 超级优先权。
                        \end{enumerate}
                  \item \zd{最高额}担保（抵押/质押/保证）。一次担保一次结算，只看结果，解决未来 n 年连续的债权债务关系。
                        \begin{enumerate}
                          \item 担保对象，持续性交易。\al{银行一次性授信}。
                          \item 不随某一具体债权的转让而转让，\zd{只随基础法律关系一同转让}。
                        \end{enumerate}
                  \item 借新还旧。新增制度。\al{银行处理陈年旧账}。
                \end{enumerate}
          \item 抵押权和租赁权的冲突。先租后押 \yjt 买卖/抵押不破租赁；\zd{先押后租 \yjt 登记破租赁}。
          \item \zd{抵押财产的转让}。\al{房贷没还完卖房}。
                \begin{enumerate}
                  \item 制度修改，同意 \yjt 通知，重大影响 \yjt 同意。
                  \item \zd{没有物上代位权}（不属于毁损、灭失、被征收）。
                \end{enumerate}
          \item 不动产登记簿内容。所有权人；他项权利人；备注，抵押合同登记。
          \item 抵押权的处分。顺位问题：
                \begin{enumerate}
                  \item 清偿顺位。不动产登记顺序；动产登记 > 未登记 > 同一顺序按比例。
                  \item 交换顺位。\zd{换位置不换钱}；要征得第三人同意。不同意不得产生不利影响。\al{丙 \yjt 乙 \yjt 丙 \yjt 甲}。
                  \item 顺位升进。
                  \item 顺位保留。购买抵押财产\zd{不影响顺位}。\al{房 85，甲 10，乙 30，丙 50，甲买走房子}。抵押权所有权混同。
                \end{enumerate}
          \item 抵押权的实现。1 个选项。\zd{和主债权的诉讼时效，同生同死}。
        \end{enumerate}
  \item 质押权。偶尔考一下。动产质押、权利质押、转质。\zd{不准占有改定}。
        \begin{enumerate}
          \item 不得对抗说 \yjt \zd{质权人将质物返还于出质人之后，不得以其质权对抗第三人}；法硕 \yjt 丧失权力说。
          \item vs 抵押。必须移转标的物的占有（交付）。
          \item 交谁处置 \yjt 孳息收取权（可以拿来抵债的意思）；问蛋归谁所有 \yjt 所有权。
          \item 质权人义务。妥善保管\footnote{实践中，一般不要动产质权，成本、风险太大。}、及时行权；出质人。\zd{催促行权}。\al{标的物价格下降}。
          \item 权利质权。有价证券（交付 \hz 登记）、可以转让的基金或股权（登记时）、应收账款\footnote{制度同保理合同、债权转让。}（债权，登记）。
          \item 转质。未征得同意 \yjt 责任转质（本质善意取得）；同意 \yjt 承诺转质（本质有权处分）。
        \end{enumerate}
  \item 留置权。唯一 1 个法定担保。1 道题目，1 个选项。
        \begin{enumerate}
          \item 只限于动产；只能留置和债权相当的金额。
          \item 民事留置。考成立要件，\zd{关键词：扣押}。
                \begin{enumerate}
                  \item 债权已到期。
                  \item 合法（\zd{有合同关系}）占有（\zd{要移转标的物占有}）债务人的动产。
                        \begin{enumerate}
                          \item 租房留置家具是不允许的。
                          \item 捡到东西返还要求先给付。
                          \item 占有、持有（占有辅助人，隶属关系）不能行使。
                          \item 合法反面：侵权行为。
                        \end{enumerate}
                  \item \zd{动产的占有和债权的发生基于同一法律关系}（民事法律关系 + 标的物）。\al{足球踢碎花瓶，扣留足球}。
                  \item 三个不，不违法、不违约、不背俗。
                  \item 可约定不准行使留置权。\al{节假日特殊物品}（身份证、人事档案、尸体、骨灰等违背公序良俗除外）。
                \end{enumerate}
          \item 商事留置。考点。
                \begin{enumerate}
                  \item 留置债务人。债权已到期；合法占有；营业债权；三个不。\zd{不需要同一法律关系}（商事要的是效率）。
                  \item 留置第三人。\zd{要基于同一法律关系}。
                \end{enumerate}
          \item 至少 60 天以上（保质期短除外，鲜活易腐）。
          \item 消灭。丧失占有；另行提供担保。
        \end{enumerate}
  \item 留置 > 购买价金 > 质押，抵押（公示在先）> 未登记动产抵押。
\end{enumerate}
